\documentclass{report}

\usepackage{amsmath,amssymb,amsthm}
\usepackage{hyperref}
\usepackage[capitalize]{cleveref}

% Theorem environments
\newtheorem{theorem}{Theorem}
\newtheorem{definition}{Definition}
\newtheorem{axiom}{Axiom}
\newtheorem{corollary}{Corollary}
\newtheorem{remark}{Remark}

% Custom macros for leanblueprint
\providecommand{\home}[1]{}
\providecommand{\github}[1]{}
\providecommand{\dochome}[1]{}
\providecommand{\lean}[1]{}
\providecommand{\leanok}{}
\providecommand{\uses}[1]{}

\title{Reverse Rosetta: Decoder Theory and Semantic Invariance in Closed and Open Systems}
\author{Fabian Schindler}
\date{\today}

\begin{document}
\maketitle

\begin{abstract}
We explore the boundaries of semantic reconstructibility and representation in formal systems. We introduce the concept of Operational Closure to mathematically distinguish closed systems (which factor uniquely through a view-independent invariant layer) from open, anchor-dependent systems. We show that while closed dialects (such as lambda calculus and set-theoretic semantics) are fully reconstructible from their invariant quotients, open systems (such as natural languages or undeciphered historical scripts like Linear A) cannot be decoded without external contextual anchors.
\end{abstract}

\chapter{Introduction and Foundational Philosophy}

Modern software engineering, programming language design, and formal verification are plagued by a profound fragmentation. Disjoint computational models—such as the untyped lambda calculus, term rewriting systems (TRS) and e-graphs, compiler intermediate representations (IR), arithmetic and quantity calculi, and virtual machine bytecodes—are treated as separate worlds. They occupy distinct namespaces, employ incompatible metatheories, and rely on ad-hoc, pairwise compiler pipelines that leak implementation details and fail to preserve observational boundaries.

This thesis argues that this fragmentation is unnecessary. We show that a broad class of closed formal systems can be represented through a single, quotient-mediated semantic kernel without loss of operational behavior. This kernel is not a new ``universal language,'' but rather a view-independent substrate: different programming languages and formalisms become different \textit{decoders} over the same invariant structure.

\section{Primal Axioms and Ontological Closure}

The ISAR (Invariant Scalable Fractal Addressable Representations) framework starts from the minimal assumptions of existence and distinction, leading to a mathematically closed computational substrate.

\begin{axiom}[Existence]
\leanok
The universe of discourse is non-empty: $\Omega \neq \emptyset$.
\end{axiom}

\begin{axiom}[Identity]
\leanok
The initial state is a singular, indistinguishable point (the Void, denoted $\emptyset$). Distinction begins when this Void is partitioned.
\end{axiom}

\begin{axiom}[Self-Description]
\leanok
A system exists only if it can describe its own existence. This self-description requirement forces the transition from a singleton to a pair: $P = (\emptyset, \emptyset)$.
\end{axiom}

\begin{axiom}[Asymmetry]
\leanok
The act of description creates a fundamental distinction between the \textit{Observer} and the \textit{Observed}. In set-theoretic terms, this corresponds to the Kuratowski pairing:
$$(x, y) = \{\{x\}, \{x, y\}\}$$
This asymmetry, denoted $\Delta$, serves as the operational fuel driving directed rewriting steps and computations.
\end{axiom}

\section{The 4-Carrier Minimality Proof by Exclusion}

The computational substrate is spanned by four primitive carriers: Identity ($I$), Rewrite ($R$), Adjacency ($A$), and State ($S$). We prove that exactly four independent carriers are necessary and sufficient to form an admissible computational substrate:

\begin{enumerate}
    \item \textbf{Identity ($C_I$)}: Represents the invariant operational equivalence. Without $C_I$, the substrate lacks a notion of stable normal forms or equivalence-preserving transformations, causing all state transitions to dissolve into noise.
    \item \textbf{Rewrite ($C_R$)}: Represents the directed, non-invertible selection/choice. Without $C_R$, transitions are symmetric, leading to static, reversible systems incapable of representing one-way computational steps.
    \item \textbf{State ($C_S$)}: Represents repeatable observables and state-space closure. Under self-application, the system must remain closed. Without $C_S$ to copy and distribute terms over contexts, the system collapses into a strictly linear, non-duplicating fragment that is sub-universal (incapable of simulating arbitrary Turing-complete computation).
    \item \textbf{Adjacency ($C_A$)}: Represents causal connectivity and topological interaction. It maps the relations between parameters. Without $C_A$, the carriers $C_I, C_R, C_S$ cannot be composed into ordered causal chains; the system collapses to an unordered bag of static values.
\end{enumerate}

No carrier can serve dual roles without violating single-carrier localization contexts:
\begin{itemize}
    \item If $C_I = C_S$, there is no proper quotient structure (observational equivalence collapses to identity).
    \item If $C_R = C_A$, causality becomes mutable, leading to non-deterministic, time-dependent topological rewrites that violate confluence.
    \item If $C_I = C_R$, the invariant becomes strategy-dependent, causing evaluation order to determine the semantic outcome.
\end{itemize}
Thus, the minimal size of the carrier set $C$ is exactly 4.

\section{Sizing Hierarchies: The Genesis of Nibbles and Bytes}

The token hierarchy of digital computing (bits, nibbles, bytes) is typically justified by legacy hardware constraints (e.g., ASCII encoding, memory alignment on the IBM System/360 in 1964). In the ISAR framework, this hierarchy is derived constructively from first principles of directed relation mapping:

\begin{itemize}
    \item \textbf{1 Bit}: Distinguishes existence from absence (atom vs. void).
    \item \textbf{2 Bits}: Identifies one of the four active roles ($C_I, C_R, C_A, C_S$) in the minimal carrier set.
    \item \textbf{4 Bits (One Nibble)}: Represents a directed pair of names (an input pointer to a carrier role and an output pointer from a carrier role). This is the minimum structure needed to express \textit{adjacency} (a relation between two roles).
    \item \textbf{8 Bits (One Byte)}: Represents a pair of pairs—a source context and a target context—which is the minimal faithful representation of a \textit{directed rewrite step} (a morphism between two carrier contexts).
\end{itemize}
Thus, early hardware engineers arrived at the 8-bit byte because it is the minimal structural representation of an asymmetric rewrite step ($\Delta$) with a parity check.

\section{Occam and the Description-Length Functional}

For any substrate $M$, we define the description-length marginal likelihood:
$$ p(D \mid M) = \int p(D \mid \theta, M) p(\theta \mid M) d\theta $$
where $\theta$ represents the parameter space (matrices and rules) and $D$ represents a reference computation. Given the 4-carrier limit, any substrate with $n > 4$ carriers introduces redundant gauge flexibility that increases description length. The terminal ISAR kernel minimizes this description length, acting as the unique Minimum Description Length (MDL) substrate.
\chapter{Decoder Theory and the Reverse Rosetta Boundary}

In this chapter, we explore the semantic boundaries of representation in formal systems. We analyze the distinction between closed computational systems—which factor through a view-independent invariant layer—and open, anchor-dependent systems.

\section{Operational Closure vs. Environmental Anchor-Dependence}

We classify computational and semantic systems into two fundamental categories based on their transition dynamics and semantic recovery:

\begin{definition}[Transition System]
\lean{ISAR.TransitionSystem}
An autonomous transition system consists of a type of states and a step relation:
$$\text{step} : \text{State} \to \text{State} \to \text{Prop}$$
\end{definition}

\begin{definition}[Operational Closure]
\lean{ISAR.IsClosed}
A subset of states $C$ is operationally closed (or forward invariant) under the transition system if all transitions starting within $C$ remain in $C$:
$$\text{IsClosed}(TS, C) \triangleq \forall (s_1, s_2 : \text{State}), C(s_1) \to \text{step}(s_1, s_2) \to C(s_2)$$
\end{definition}

We prove that states starting in an operationally closed subsystem remain within it under any sequence of transitions:

\begin{theorem}[Forward Invariance of Closed Subsystems]
\lean{ISAR.closure_preserved_under_reachability}
\uses{ISAR.IsClosed}
For any transition system $TS$ and closed subsystem $C$, any state reachable from a state in $C$ remains in $C$:
$$\forall (s_1, s_2 : \text{State}), C(s_1) \to \text{Reachable}(TS, s_1, s_2) \to C(s_2)$$
\end{theorem}

Operational closure forms the basis for reconstructibility. In a closed system (such as lambda calculus, SKI combinators, and set-theoretic semantics), the transition dynamics can be fully projected onto the $\text{InvariantLayer}$ quotient, allowing the decoder to reconstruct semantics at any point without loss of meaning.

In contrast, open systems depend on external parameters:

\begin{definition}[Anchor-Dependent System]
\lean{ISAR.AnchorDependentSystem}
A transition system where transitions depend on an external environment or context:
$$\text{step} : \text{State} \to \text{Anchor} \to \text{State} \to \text{Prop}$$
\end{definition}

In an anchor-dependent system, trace semantics under sequences of external anchors are non-deterministic from the perspective of the state alone:

\begin{theorem}[Referential Openness Requires Anchors]
\lean{ISAR.referentially_open_requires_anchor}
If a state can lead to different observations under different anchor sequences, then the trace semantics are not recoverable from the state alone:
$$\exists (\text{as} : \text{List Anchor}), \neg \left(\forall (s' : \text{State}), \text{Trace}(s, \text{as}, s') \to \text{decode}(s') = \text{decode}(s_1')\right)$$
\end{theorem}

\section{The Boundary of Decodability (Reverse Rosetta)}

The **Reverse Rosetta** principle represents the boundary of what can be decoded when only the syntax is preserved but the semantic context is lost.

For closed systems (like lambda calculus or hereditarily finite sets), we can reconstruct their semantics from the operational quotient layer because their reductions are self-contained. For open systems (such as natural languages, undeciphered scripts like Linear A, or APIs with external state), the characters and syntax do not suffice for decoding because the meaning is tied to external cultural or environmental anchors.

\section{Formal Distinction between Invariant Layer and Decoders}

The distinction between the **Invariant Layer** and the **Decoder** is formalized through three distinct mechanisms:

\subsection{Category-Theoretic Terminality}
As proved in Theorem 9, the quotient $\text{InvariantLayer}$ is the terminal object in the category of semantic kernels. Any other admissible computational view is guaranteed to factor uniquely through it via a unique morphism:
$$h : K \to \text{ISAR\_Kernel}$$

\subsection{Geometrical Gauge Invariance}
As formalized in the matrix representation of carriers, the decoder corresponds to a coordinate system (or basis choice), while the Invariant Layer represents the coordinate-free operator. Two matrix representations are isomorphic under basis conjugation:
$$P \cdot K_1 \cdot P^{-1} = K_2$$
proving that representation details are gauge-dependent shadows.

\subsection{Empirical Decoupling}
The type $\text{InvariantLayer}$ remains static, yet we implement entirely separate decoders for set membership, lambda terms, and stack-machine bytecodes. These map disjoint syntactic types into the same quotient, proving that the invariant layer holds the view-independent semantics while representation remains strictly delegated to the decoder.

\chapter{Discussion, Related Work, and Limitations}

We present a review of related literature, establish the formal boundaries of validity for our theorems, and outline directions for future research.

\section{Related Work}

The ISAR framework intersects several foundational areas of computer science and logic:

\begin{itemize}
    \item \textbf{Combinatory Logic and Lambda Calculus}: The bracket abstraction algorithm compiled to combinator form is a classical approach pioneered by Curry and Feys~\cite{curry1958combinatory} and Barendregt~\cite{barendregt1984lambda}. While traditional implementations suffer from code size explosion (e.g., $O(n^3)$ for standard abstraction), the linear fragment of ISAR maintains bounded sizing, sharing a direct connection to optimal reduction.
    \item \textbf{Optimal Evaluation and Interaction Nets}: Yves Lafont's interaction nets~\cite{lafont1990interaction} provide a graphical formulation of computation where reductions are local and symmetric. Our formalization of the linear fragment (\texttt{LinearIKTerm}) and its bounds mirrors the behavior of optimal virtual machines like the Higher-Order Virtual Machine 2 (HVM2)~\cite{hvm2}, which compiles terms to interaction nets for parallel execution.
    \item \textbf{Categorical Semantics of Computation}: Lawvere's functorial semantics~\cite{lawvere1963functorial} established algebraic theories as categories. Our formulation of semantic kernels ($\mathcal{K}$) and the proof of terminality of the Invariant Layer quotient formalizes view-independence as a universal factorization property.
    \item \textbf{Proof Assistant Kernels}: The verification and automated generation of documentation utilizing Lean 4 quotients and metadata mappings (e.g., doc-gen4~\cite{docgen4}) represent the state-of-the-art in proof design. Our work extends this by proving formal simulation theorems directly against the proof assistant's environment.
\end{itemize}

\section{Limitations and Scope Boundaries}

The theorems presented in this work are accompanied by strict boundary conditions:

\begin{enumerate}
    \item \textbf{Closed Dialects Only}: The terminality and factorization theorems apply exclusively to computational dialects that are \textit{operationally closed}—meaning they possess a total encoding function into the substrate and their evaluation dynamics can be modeled as deterministic, confluent rewrite steps. This excludes open, dynamic, or referentially open systems (such as natural languages, interactive APIs, or non-deterministic transition systems) that require contextual anchors.
    \item \textbf{ZFC HF Set Scope}: The set-theoretic relative interpretation is verified solely for the Hereditarily Finite set fragment. Extension to infinite sets (ZFC with the axiom of infinity) requires different ordinals and is outside the scope of the current Lean 4 proofs.
    \item \textbf{Interpretive Physical Assumptions}: The application of the Universal Approximation Theorem to physical systems rests on the interpretive hypothesis that a physical system's state space can be modeled as an admissible continuous kernel. This hypothesis is a modeling choice and is not mechanically verified by the Lean 4 proof assistant.
\end{enumerate}

\section{Future Work}

This paper forms the core substrate of a multi-part series of monographs:
\begin{itemize}
    \item \textbf{Paper 2 (Continuous Approximation)}: Extension of the discrete kernel to real-valued physical quantities, continuous address spaces, and metric completions, proving the unique continuous representation limit. This builds upon classical universal approximation results on compact domains, particularly those of Cybenko~\cite{cybenko1989approximation} and Leshno et al.~\cite{leshno1993multilayer}.
    \item \textbf{Paper 3 (Hardware Compilation)}: Verification of the compilation path from the algebraic matrix representations to optimal virtual machines (HVM2) and relational database schemas. We also formalize partial evaluation and dialect specialization via the first, second, and third Futamura projections~\cite{futamura1971partial}.
\end{itemize}

\bibliographystyle{plain}
\bibliography{references}



\end{document}
